% from ega4-18-fr.tex
\oldpage[IV-4]{165}

on en déduit que $X'$ est réduit (8.7.1). Cela étant, si $n(y)=n$, on a $n_s=n$ et par suite les corps résiduels aux points de $X_\eta$ sont nécessairement séparables sur $k(\eta)$; les corps résiduels aux points de $X'_{\eta'}$ sont par suite algébriques et séparables sur $k(\eta')$ (4.6.1), autrement dit $n'_s=n'=n$, et on a donc en vertu de ce qui précède $n'(y')=n'$. Le morphisme $f'$ a donc les mêmes propriétés que $f$, et si l'on prouve que $f'$ est étale, il en résultera que $f$ est étale par (17.7.1).

\emph{D) Fin de la démonstration.} --- Supposons donc $A$ strictement local. Puisque $A$ est hensélien et le morphisme $f$ fini, il résulte de (18.5.11) que si $x_j$ ($1\leq j\leq m$) sont les points distincts de $f^{-1}(y)$, $X$ est somme des sous-préschémas ouverts $\Spec(\sh{O}_{X,x_j})=X_j$, qui sont donc les composantes connexes de $X$. Par hypothèse, chacune d'elles domine $Y$, donc rencontre $f^{-1}(\eta)$, et par suite le nombre de points de $f^{-1}(\eta)$ est $\geq m$; \emph{a fortiori} on a $n_s\geq m$. Mais puisque $k(y)$ est séparablement clos et que les $k(x_j)$ sont algébriques sur $k(y)$, ce sont nécessairement des extensions radicielles de $k(y)$, donc $m=n(y)$, et cela prouve (i). Supposons maintenant vérifiées les hypothèses de (ii); les relations $n(y)=n_s=n$ entraînent d'abord, puisque toute composante irréductible de $X$ domine $Y$, que chacun des $X_j$ est irréductible; en outre, comme les $X_j$ sont réduits, ils sont intègres; enfin les relations $m=n_s=n$ entraînent que pour le point générique $\xi_j$ de $X_j$, $k(\xi_j)$ est une extension séparable de $k(\eta)$ de degré 1, donc isomorphe à $k(\eta)$; autrement dit, la restriction $f|X_j:X_j\to Y$ de $f$ est un morphisme fini et birationnel; en outre, $X_j$ est intègre et par hypothèse $Y$ est normal; donc (8.12.10.1) $f|X_j$ est un isomorphisme, ce qui prouve que $f$ est étale, et termine la démonstration de (18.10.16).

\begin{remark}[18.10.17]
\label{IV.18.10.17}
La méthode de «localisation étale» utilisée dans la démonstration de (18.10.16) permet aussi d'améliorer les résultats de (15.5.1) en éliminant l'hypothèse noethérienne. On a tout d'abord le résultat suivant qui généralise (15.5.2):
\end{remark}

\begin{env}[18.10.17.1]
\label{IV.18.10.17.1}
Soient $f:X\to Y$ un morphisme séparé et quasi-fini, $y$ un point de $Y$ tel que $f$ soit universellement ouvert aux points de $f^{-1}(y)$. Alors, pour toute générisation $y'$ de $y$, on a $n(y)\leq n(y')$.
\end{env}

On peut remplacer $Y$ par le sous-préschéma fermé réduit $Z$ ayant pour espace sous-jacent $\overline{\{y'\}}$ et $X$ par $f^{-1}(Z)$, donc supposer $Y$ intègre de point générique $y'$. Utilisant (18.12.13) comme dans (18.10.16, B)), on peut supposer que $X$ est ouvert dense dans un préschéma $X'$ et que $f$ est la restriction à $X$ d'un morphisme fini $g:X'\to Y$. Montrons que l'on a $g^{-1}(y')=f^{-1}(y')$. En effet, soit $i:g^{-1}(y')\to X'$ l'injection canonique, qui est un morphisme plat puisque $y'$ est point générique de $Y$ (I, 3.6.5); on peut écrire $f^{-1}(y')=i^{-1}(X)$. D'autre part, l'injection canonique $j:X\to X'$ est un morphisme de type fini puisque le composé $g\circ j=f$ est de type fini et que $g$ est séparé (1.5.4); donc $X$ est un ouvert rétrocompact dans $X'$, et par suite est pro-constructible dans $X'$ (1.9.5, (v)). On conclut donc de (2.3.10) que l'on a $i^{-1}(\overline{X})=\overline{i^{-1}(X)}$, autrement dit $f^{-1}(y')$ est dense dans $g^{-1}(y')$; mais comme $g^{-1}(y')$ est discret, on a nécessairement $f^{-1}(y')=g^{-1}(y')$.

On est ainsi ramené à prouver que la somme des nombres $[k(x):k(y)]_s$ où $x$ parcourt l'ensemble des points de $g^{-1}(y)$ où $g$ est universellement ouvert, est au plus égale

